\documentclass{../../slides-style}

\slidetitle{Практика по графам}{07.03.2026}

\begin{document}
    
    \begin{frame}[plain]
        \titlepage
    \end{frame}
    
    \begin{frame}{Задача на пару}
        Реализовать обход графа в ширину
        \begin{outline}
            \1 Граф задан в файле виде матрицы смежности
            \1 Как параметр (через argc/argv) читается стартовая вершина
            \1 Задача --- распечатать вершины графа в порядке обхода в ширину (в каком-то из порядков)
            \1 Нужен АТД \enquote{Граф} (и АТД \enquote{Очередь})
            \1 Работаем всей группой, по очереди выходим к проектору и вносим свой вклад
                \2 За бонусный балл к домашке
        \end{outline}
    \end{frame}

\end{document}